% Source: physical PDF pp. 360--364 (printed pp. 337--341), beginning at
% the Section 31.4 heading on p. 360 and ending immediately before the
% Section 31.5 heading on p. 364.
% Inventory: Eqs. (31.4.1)--(31.4.20), three linked citation-marker
% occurrences (references [11a], [11b], and [11b]), and no unnumbered
% displays, footnotes, figures, tables, or centered three-asterisk dividers.
% Source wording preserved: the opening paragraph says ``minimum
% supersymmetric standard model.'' Direct products of gauge groups use
% \times, as in the scan; \cdot is reserved for ordinary or scalar
% multiplication and inner products.
% Uncertainties: none.

\subsection{Anomaly-Mediated Supersymmetry Breaking}

In Section 28.3 the possibility was raised that supersymmetry may be broken
in some sort of hidden sector of superfields that do not carry the
\(SU(3)\times SU(2)\times U(1)\) quantum numbers of the
standard model, and communicated to observable particles gravitationally.
In this section we will deal with one class of supersymmetry-breaking
effects in the minimal supersymmetric standard model, those of first order
in \(\kappa\equiv\sqrt{8\pi G}\). This includes the gaugino masses and the
parameters \(A_{ij}\) and \(B\) in the Lagrangian density
\eqref{eq:28.4.1}. Other supersymmetry-breaking effects such as squark and
slepton squared masses are of second order in \(\kappa\), and will be taken
up in Section 31.7, when we consider gravity-mediated supersymmetry breaking
using the general supergravity formalism described in Section 31.6.

We can find the effects of gravity-mediated supersymmetry breaking to first
order in \(\kappa\) by simply replacing the component fields of the
gravitational supermultiplet in the interaction \eqref{eq:31.1.34} with
their expectation values. The only ones of these component fields that can
acquire non-vanishing vacuum expectation values from the spontaneous
breakdown of supersymmetry in a hidden sector of matter superfields are the
spinless auxiliary fields \(s\) and \(p\), so with Eq.~\eqref{eq:31.1.33}
this gives a first-order supersymmetry-breaking interaction
\begin{equation}
  \mathcal L^{(1)}
  =2\kappa\bigl[-A^X\langle p\rangle+B^X\langle s\rangle\bigr]
  =3\operatorname{Im}
   \bigl[\widetilde m_g^*(A^X+iB^X)\bigr],
  \tag{31.4.1}\label{eq:31.4.1}
\end{equation}
where \(\widetilde m_g\) is the complex gravitino mass
\eqref{eq:31.3.20}, and \(A^X\) and \(B^X\) are the \(A\)- and
\(B\)-components of the real chiral scale-non-invariance superfield \(X\)
discussed in Section 26.7.

We showed in Section 26.7 that, for a renormalizable theory of left-chiral
superfields \(\Phi_n\) with superpotential \(f(\Phi)\), the \(X\) superfield
is given by
\begin{equation}
  X=\frac23\operatorname{Im}
  \left[
    \sum_n\Phi_n\frac{\partial f(\Phi)}{\partial\Phi_n}
    -3f(\Phi)
  \right].
  \tag{31.4.2}\label{eq:31.4.2}
\end{equation}
This can be put in a form that allows an immediate extension to more
general theories by writing the coupling parameters of the superpotential
in terms of dimensionless parameters and a parameter \(\mathcal M\) with
the dimensions of mass. Since the superpotential has the dimensions of
\((\mathrm{mass})^3\), we have
\begin{equation}
  \mathcal M\frac{\partial f(\Phi)}{\partial\mathcal M}
  +\sum_n\Phi_n\frac{\partial f(\Phi)}{\partial\Phi_n}
  =3f(\Phi).
  \tag{31.4.3}\label{eq:31.4.3}
\end{equation}
Therefore Eq.~\eqref{eq:31.4.2} may be written
\begin{equation}
  X=\frac23\operatorname{Im}
  \left[
    \mathcal M\frac{\partial f(\Phi)}{\partial\mathcal M}
  \right].
  \tag{31.4.4}\label{eq:31.4.4}
\end{equation}
This formula can be generalized to give the contribution to \(X\) of the
scale dependence of any sort of \(\mathcal F\)-term in the Lagrangian. A
term \(2\operatorname{Re}[f(\Phi,W)]_{\mathcal F}\) in the Lagrangian makes
a contribution to \(X\) given by the obvious generalization of
Eq.~\eqref{eq:31.4.4}:
\begin{equation}
  X=\frac23\operatorname{Im}
  \left[
    \mathcal M\frac{\partial f(\Phi,W)}{\partial\mathcal M}
  \right].
  \tag{31.4.5}\label{eq:31.4.5}
\end{equation}
(There is also a contribution to \(X\) from any mass scale dependence in
the \(D\)-terms in the Lagrangian.) Comparing Eq.~\eqref{eq:31.4.5} with
Eqs.~\eqref{eq:26.3.10} and \eqref{eq:26.3.13}, we see that
\begin{equation}
  A^X+iB^X
  =\frac{2\mathcal M}{3i}\frac{\partial}{\partial\mathcal M}
   [f(\Phi,W)]_{\theta=0}
  =\frac{2\mathcal M}{3i}
   \frac{\partial f(\phi,\lambda_L)}{\partial\mathcal M}.
  \tag{31.4.6}\label{eq:31.4.6}
\end{equation}
There will then be a supersymmetry-breaking term in the effective
Lagrangian of first order in \(\kappa\), given by
Eqs.~\eqref{eq:31.4.1} and \eqref{eq:31.4.6} as
\begin{equation}
  \mathcal L_f^{(1)}
  =-2\operatorname{Re}
  \left[
    \widetilde m_g^*
    \frac{\mathcal M\,\partial f(\phi,\lambda_L)}
         {\partial\mathcal M}
  \right].
  \tag{31.4.7}\label{eq:31.4.7}
\end{equation}

Eq.~\eqref{eq:31.4.7} applies not only for terms in the Lagrangian with
explicit scale dependence, but also for the scale dependence of coupling
constants described by the renormalization
group.~\hyperref[ch31-ref-11a]{[11a]} This scale dependence arises from the
quantum mechanical anomaly that gives a non-zero value to the trace of the
energy-momentum tensor and also to the divergence of the \(R\)-current, so
the observable supersymmetry-breaking effects that arise in this way are
said to be \emph{anomaly-mediated}.

Consider in particular the kinematic term \(\mathcal L_{\mathrm{gauge}}\)
for renormalizable supersymmetric gauge theories, given by
Eqs.~\eqref{eq:27.3.22} and \eqref{eq:27.3.23} as
\begin{equation}
  \mathcal L_{\mathrm{gauge}}
  =-\frac1{2g^2}\operatorname{Re}
  \sum_{A\alpha\beta}
  [\epsilon_{\alpha\beta}W_{A\alpha}W_{A\beta}]_{\mathcal F}.
  \tag{31.4.8}\label{eq:31.4.8}
\end{equation}
This does not depend explicitly on any mass scale, but the coupling
constant \(g\) has the familiar dependence on a renormalization scale
\(\mathcal M\), given by the renormalization group equation
\begin{equation}
  \mathcal M\frac{dg(\mathcal M)}{d\mathcal M}
  =\beta\bigl(g(\mathcal M)\bigr).
  \tag{31.4.9}\label{eq:31.4.9}
\end{equation}
Then Eq.~\eqref{eq:31.4.7} shows that the gauge Lagrangian
\eqref{eq:31.4.8} yields a supersymmetry-breaking term in the Lagrangian
\begin{equation}
  \mathcal L_{\mathrm{gauge}}^{(1)}
  =-\frac{\beta(g)}{g^3}\operatorname{Re}
  \left[
    \widetilde m_g^*
    \sum_{A\alpha\beta}
    \epsilon_{\alpha\beta}\lambda_{AL\alpha}\lambda_{AL\beta}
  \right],
  \tag{31.4.10}\label{eq:31.4.10}
\end{equation}
where \(\lambda_{AL\alpha}\) is the left-handed part of the gaugino field,
normalized like \(W_{A\alpha}\) by multiplication with the gauge coupling
\(g\), so that \(g\) does not appear in the structure constants or the
interaction of the gauge superfield with quark superfields. Taking account
of this normalization convention, the gaugino mass equals \(g^2\) times
the absolute value of the coefficient of
\(\frac12\sum_{A\alpha\beta}
\epsilon_{\alpha\beta}\lambda_{AL\alpha}\lambda_{AL\beta}\), or
\begin{equation}
  m_{\mathrm{gaugino}}
  =m_g\left|\frac{\beta(g)}{g}\right|.
  \tag{31.4.11}\label{eq:31.4.11}
\end{equation}
In this formula \(m_{\mathrm{gaugino}}\) and \(g\) are cut-off-dependent
bare parameters in the Lagrangian, governed by Wilsonian renormalization
group equations. In Section 27.6 we saw that \(\beta(g)\) arises purely
from one-loop diagrams, so that \(\beta(g)=bg^3\), with \(b\) a constant,
and therefore
\begin{equation}
  m_{\mathrm{gaugino}}=m_g|b|g^2.
  \tag{31.4.12}\label{eq:31.4.12}
\end{equation}
The physical mass of the gaugino differs from Eq.~\eqref{eq:31.4.12} by
corrections of higher order in \(g\) but, since we know that gauginos must
be much heavier than the characteristic scale of quantum chromodynamics,
these corrections are small for the gluino as well as for the wino and
bino.

\begingroup
\emergencystretch=1em
Supersymmetry does not allow any explicitly scale-dependent terms in the
Lagrangian for the gluon and quark superfields, so, with electroweak
interactions neglected, Eq.~\eqref{eq:31.4.12} gives the only contribution
of order \(\kappa\) to the gluino mass. For three quark generations,
Eq.~\eqref{eq:28.2.10} gives \(b=-3g_s^3/16\pi^2\). Taking
\(g_s^2/4\pi=0.118\), Eq.~\eqref{eq:31.4.12} gives the gluino mass as
\par
\endgroup
\begin{equation}
  m_{\mathrm{gluino}}
  =\frac{3g_s^2m_g}{16\pi^2}
  =2.8\times10^{-2}m_g.
  \tag{31.4.13}\label{eq:31.4.13}
\end{equation}
On the other hand, there is a scale-dependent interaction in the
Lagrangian for the Higgs superfields, produced by the \(\mu\)-term
\(-\mu(\mathcal H_2^{\mathsf T}e\mathcal H_1)\) in
Eq.~\eqref{eq:28.1.7}, so there is a term in the Lagrangian given by
Eq.~\eqref{eq:31.4.7} as
\(2\operatorname{Re}
[\widetilde m_g^*\mu(\mathcal H_2^{\mathsf T}e\mathcal H_1)]\).
Comparing with the \(B\mu\)-term in Eq.~\eqref{eq:28.4.1}, we see that
this gives
\begin{equation}
  B=-\widetilde m_g^*.
  \tag{31.4.14}\label{eq:31.4.14}
\end{equation}

With three generations of quarks and lepton superfields and one pair of
Higgs doublet superfields \(H_1\) and \(H_2\),
Eqs.~\eqref{eq:28.2.8} and \eqref{eq:28.2.9} give \(b=11/16\pi^2\) for
the \(U(1)\) gauge coupling \(g'\) and \(b=1/16\pi^2\) for the \(SU(2)\)
gauge coupling \(g\). Taking \(g'^2/4\pi=0.0102\) and
\(g^2/4\pi=0.0338\), Eq.~\eqref{eq:31.4.11} would give masses
\(8.9\times10^{-3}m_g\) and \(2.7\times10^{-3}m_g\) for the bino and
wino, respectively. However, the bino and wino masses also receive
contributions of order \(g'^2m_g/16\pi^2\) and \(g^2m_g/16\pi^2\),
respectively, from diagrams in which the bino or wino is attached to a
higgs--higgsino loop, with supersymmetry breaking introduced by the term
\(2\operatorname{Re}
[\widetilde m_g^*\mu(\mathcal H_2^{\mathsf T}e\mathcal H_1)]\)
in the Lagrangian density. This gives the bino and wino
masses~\hyperref[ch31-ref-11b]{[11b]}
\begin{equation}
  m_{\mathrm{bino}}
  =\frac{g'^2m_g}{16\pi^2}
   \left|11-f\left(\frac{\mu^2}{m_A^2}\right)\right|,
  \tag{31.4.15}\label{eq:31.4.15}
\end{equation}
\begin{equation}
  m_{\mathrm{wino}}
  =\frac{g^2m_g}{16\pi^2}
   \left|1-f\left(\frac{\mu^2}{m_A^2}\right)\right|,
  \tag{31.4.16}\label{eq:31.4.16}
\end{equation}
where \(m_A\) is the pseudoscalar particle mass defined by
Eq.~\eqref{eq:28.5.21}, and
\begin{equation}
  f(x)\equiv\frac{2x\ln x}{x-1}.
  \tag{31.4.17}\label{eq:31.4.17}
\end{equation}
The implications of these results will be considered further in Section
31.7.

Finally, there is a scale-dependent field-renormalization factor \(Z_r\)
multiplying the kinematic Lagrangian density
\([\Phi_r^*e^{-V}\Phi_r]_D\) for any left-chiral superfield \(\Phi_r\).
This can be moved from the kinematic \(D\)-terms to the superpotential
\(\mathcal F\)-terms by absorbing a factor \(Z_r^{1/2}\) in \(\Phi_r\).
The Yukawa couplings \(h_{rst}\) (such as the \(h_{ij}^E\), \(h_{ij}^D\),
and \(h_{ij}^U\) in Eq.~\eqref{eq:28.1.7}) in the trilinear
superpotential terms \(\sum_{rst}h_{rst}\Phi_r\Phi_s\Phi_t\) are then
multiplied with a factor \(Z_r^{-1/2}Z_s^{-1/2}Z_t^{-1/2}\) that depends
on the cut-off \(\mathcal M\). According to Eq.~\eqref{eq:31.4.7}, the
interaction \(\mathcal L^{(1)}\) will then make a contribution to the
Lagrangian density:
\begin{equation}
  \mathcal L_{\mathrm{Yukawa}}^{(1)}
  =-2\sum_{rst}\gamma_{rst}\operatorname{Re}
  [\widetilde m_g^*h_{rst}\phi_r\phi_s\phi_t],
  \tag{31.4.18}\label{eq:31.4.18}
\end{equation}
where
\begin{equation}
  \begin{aligned}
  \gamma_{rst}
  \equiv{}&
  \mathcal M\frac{\partial\ln h_{rst}(\mathcal M)}
                   {\partial\mathcal M}
  \\
  ={}&
  -\frac12\mathcal M\frac{\partial\ln Z_r(\mathcal M)}
                            {\partial\mathcal M}
  -\frac12\mathcal M\frac{\partial\ln Z_s(\mathcal M)}
                            {\partial\mathcal M}
  -\frac12\mathcal M\frac{\partial\ln Z_t(\mathcal M)}
                            {\partial\mathcal M}.
  \end{aligned}
  \tag{31.4.19}\label{eq:31.4.19}
\end{equation}
We see that the coefficients \(A_{ij}^E\), \(A_{ij}^D\), and \(A_{ij}^U\)
in Eq.~\eqref{eq:28.1.7} are given
by~\hyperref[ch31-ref-11b]{[11b]}
\begin{equation}
  A_{ij}^N
  =\widetilde m_g^*\gamma_{ij}^N
  =\widetilde m_g^*\mathcal M
   \frac{\partial\ln h_{ij}^N}{\partial\mathcal M},
  \tag{31.4.20}\label{eq:31.4.20}
\end{equation}
where \(N=E\), \(D\), or \(U\). This gives \(A_{ij}^D\) and \(A_{ij}^U\)
of order \(g_s^2m_g/8\pi^2\), while \(A_{ij}^E\) is of order
\(g^2m_g/8\pi^2\) or \(g'^2m_g/8\pi^2\).
